The present study sits at the intersection of four lines of work. The first compares human and AI emotion judgments and motivates the use of continuous scales rather than coarse categorical labels. Our earlier VAS-versus-Likert study showed that a VAS-based design is practical for short-story affect evaluation and provides a stronger foundation for fine-grained human--LLM comparison than a small discrete scale \cite{shirahama2025vas}.

The second line concerns benchmark design for LLM emotional intelligence. Benchmark suites such as EQ-Bench, EmotionBench, EmoBench, and EmotionQueen have helped standardize evaluation of emotional reasoning, empathy, and affective response quality \cite{paech2024eqbench,huang2024emotionbench,sabour2024emobench,chen2024emotionqueen}. However, these suites are not primarily designed to ask which current models are closest to a human reference profile on a literary reading task with continuous story-by-emotion scoring.

The third line concerns multilingual and robustness-oriented LLM evaluation. Language changes can alter lexical choice, narrative framing, and calibration, so multilingual comparison is necessary when a benchmark is meant to support deployment across translated conditions. Yet multilingual variation alone is not equivalent to human validity unless matched human references are also available.

The fourth line is our own prior platform-development and mathematical-characterization work. Those studies established that LLMs differ substantially in reliability, controllability, and emotional individuality across a broad model landscape \cite{openrouterplatform,shirahama2026mdpi}. The present paper intentionally narrows the design: one neutral-reader primary condition, one human-grounded language endpoint, and multilingual runs interpreted as robustness analyses rather than as direct substitutes for human validation.